\chapter{The Approximation Theorem and Its Relation to Equivalent MLC Statements}
\label{ch:approximation}
\noindent This chapter constructs maps that fix increasingly deep finite outer approximations to the Mandelbrot set and converge to it in Hausdorff distance. This yields approximations at every scale and a suitable sequence without assuming MLC. It then proves a stronger equivalence: equicontinuous approximation with a common modulus of continuity holds if and only if MLC holds. The proof uses the external conformal map, retractions onto Green sublevel sets, and compactness of families of maps.

\section{Statement of the Approximation Theorem}
\label{sec:approximation-statement}

\noindent\emph{Role of the result.} The theorem is retained as an
independent result about finite outer approximations: it establishes
\(\mathsf{Fin}\) and \(\mathsf{EA}_{\mathrm{seq}}\) in
\eqref{eq:main-discrepancy}. The constructed maps \(T_n\) are not used
in the \(\mathsf{AM}\) or \(\mathsf{GC}\) reductions, nor in the
conditional implication from \(\mathsf H_{\M}\) to factorization; those
reductions use only the dyadic grid constructed in this chapter.

Let \(\mathbf{CompHaus}\) denote the category of compact Hausdorff
spaces and continuous maps. For compact Hausdorff spaces \(X,Y\), set
\[
 \operatorname{Hom}(X,Y):=C(X,Y).
\]
If \(Y\) is metrized, equip the morphism space with the uniform metric.
Let \(\mathbf N\) be the category associated with the partially ordered
set of natural numbers. The inclusions \(\iota_N:O_N\hookrightarrow\D\)
define an inverse system
\(\mathcal O:\mathbf N^{\mathrm{op}}\to\mathbf{CompHaus}\); Step~1
identifies its limit with \(\M\).

\begin{approxthm}[A Sequence of Endomorphisms Fixing Successive Levels]
There exist \(\lambda=(\lambda_n)_{n\in\N}\in\N^\N\) and continuous
maps \(T_n:\D\to\D\), \(n\in\N\), such that
\begin{gather}
 \lambda_0=0,\qquad\lambda_1=1,\qquad
 \forall n\in\N,\quad \lambda_n\ge n,\quad
 \lambda_{n+1}>\lambda_n,\label{eq:opening-depths}\\
 T_0=\id_{\D},\qquad T_1=\Pi|_{\D},\qquad
 \forall n\ge1,\quad T_n(\D)\subseteq K,\label{eq:opening-maps}\\
 \forall n\in\N,\quad
 T_n\circ\iota_{\lambda_n}=\iota_{\lambda_n},
 \label{eq:opening-fixing}\\
 \forall n\in\N,\quad
 \sup_{z\in\D}\dist(T_n(z),O_{\lambda_n})
                 \le2^{-n},
 \qquad
 \forall n\ge1,\quad T_n\in\mathcal R_n,
 \label{eq:opening-proximity}\\
 \forall n\in\N,\quad
 \M\subseteq T_n(\D),\qquad
 \forall n\ge2,\quad
 \dH(T_n(\D),\M)\le2^{-n}.
 \label{eq:opening-image-distance}
\end{gather}
\end{approxthm}

\section{Proof of the Approximation Theorem}
\label{sec:approximation-proof}

\subsection{Step 1. Realize Finite Orbit Constraints as Compact Objects}

\begin{lemma}[Description of \(\M\) by the Sets \(O_N\)]
\label{lem:outer-sets}
\begin{gather}
 \M=\{c\in\C:\forall n\in\N,\ |p_n(c)|\le2\},\label{eq:escape}\\
 \M\subseteq O_{N+1}\subseteq O_N\subseteq K\subseteq\D,\qquad
 \bigcap_{N\in\N}O_N=\M,\label{eq:outer-nesting}\\
 \varprojlim_{\mathbf N^{\mathrm{op}}}\mathcal O
       \cong\M.\label{eq:outer-limit}
\end{gather}
The sets \(\M,O_N,K,\D\) are compact and closed, and \(\D\) is convex.
The transition morphisms in \(\mathcal O\) are the inclusions
\(\iota_{N+1,N}:O_{N+1}\hookrightarrow O_N\).
\end{lemma}
\begin{proof}
Suppose \(q:=|p_n(c)|>2\) and \(|c|\le q\). Induction gives
\[
 |p_{n+j}(c)|\ge q(q-1)^j\qquad(j\in\N).
\]
For the induction step, if \(r:=|p_{n+j}(c)|\ge q\), then
\[
 |p_{n+j+1}(c)|
 \ge r^2-|c|
 \ge r^2-r
 \ge(q-1)r.
\]
Thus the orbit is unbounded. If \(|c|>2\), take \(n=1\) and
\(q=|p_1(c)|=|c|\). If \(|c|\le2\) and some orbit iterate has modulus
greater than \(2\), use it as \(q\). This proves one direction of
\eqref{eq:escape}; the converse follows from boundedness with \(B=2\).

Each \(p_j\) is continuous, so \(p_j^{-1}(K)\) is closed. Hence \(O_N\)
is closed in the compact set \(K\). Since \(p_0(c)=0\) and \(p_1(c)=c\),
we have \(O_{N+1}\subseteq O_N\subseteq K\). Equation~\eqref{eq:escape}
gives \(\M\subseteq O_N\) and
\[
 \bigcap_{N\in\N}O_N
 =K\cap\bigcap_{j\in\N}p_j^{-1}(K)=\M.
\]
The same equality shows that \(\M\) is closed in \(K\), hence compact.
Finally, \(K\subseteq\D\), where \(\D\) is a closed, bounded, convex
square. Let \(X\) be a compact Hausdorff space and let
\(u_N:X\to O_N\) be maps satisfying
\(\forall N,\ \iota_{N+1,N}\circ u_{N+1}=u_N\). Then the composites
\(\iota_N\circ u_N:X\to\D\) all agree. Their common image lies in
\(\bigcap_NO_N=\M\), so they factor uniquely through the inclusion
\(\M\hookrightarrow\D\). This is the universal property in
\(\mathbf{CompHaus}\) asserted in~\eqref{eq:outer-limit}.
\end{proof}

\subsection{Step 2. Prove Convergence in the Space of Compact Subsets}

\begin{definition}[Hausdorff Distance to \(\M\)]
\label{def:outer-distance}
Let \(\mathcal K(K)\) be the space of nonempty compact subsets of \(K\)
with the Hausdorff metric. For \(N\in\N\), set
\[
 h_N:=\dH(O_N,\M)
     =\sup_{z\in O_N}\dist(z,\M).
\]
The sequence \((O_N)\) converges to \(\M\) in \(\mathcal K(K)\).
\end{definition}

\begin{theorem}[Convergence in Hausdorff Distance]\label{thm:outer-distance}
\[
 0\le h_{N+1}\le h_N,\qquad \lim_{N\to\infty}h_N=0.
\]
Define
\begin{gather}
 \ell_0:=0,\qquad\ell_1:=1,\qquad
 \ell_n:=\min\{N\in\N:N>\ell_{n-1},\ N\ge n,\ h_N\le2^{-n}\}
 \quad(n\ge2),\label{eq:lower-depths}\\
 \Lambda:=
 \{\lambda\in\N^\N:
   \lambda_0=0,\ \lambda_1=1,\
   \forall n,\ \lambda_n\ge\ell_n,\
   \lambda_{n+1}>\lambda_n\}.
 \label{eq:depth-class}
\end{gather}
Then every minimum in \eqref{eq:lower-depths} exists,
\(\ell\in\Lambda\), and
\[
 \forall\lambda\in\Lambda\ \forall n\ge2,\qquad
 \lambda_n\ge n,\quad h_{\lambda_n}\le2^{-n}.
\]
For every \(\lambda\in\Lambda\), we have \(\lambda_n\to\infty\) and
\(\bigcap_{n\in\N}O_{\lambda_n}=\M\). Hence the limit of the reindexed
inverse system \(\mathcal O\circ\lambda^{\mathrm{op}}\) is \(\M\).
\end{theorem}
\begin{proof}
Nestedness gives \(h_{N+1}\le h_N\). Let the limit be \(h_\infty\).
For each \(N\), choose \(z_N\in O_N\) such that
\(\dist(z_N,\M)=h_N\). If \(h_\infty>0\), compactness of \(O_0\)
gives a subsequence \(z_{N_k}\to z_\infty\in O_0\). For each fixed
\(m\), we have \(N_k\ge m\) for all sufficiently large \(k\), hence
\(z_{N_k}\in O_m\). Since \(O_m\) is closed,
\(z_\infty\in\bigcap_mO_m=\M\). Continuity of distance to a nonempty
set then gives
\[
 h_{N_k}=\dist(z_{N_k},\M)\longrightarrow0,
\]
contradicting \(h_{N_k}\to h_\infty>0\). Therefore \(h_N\to0\), and
the required integers \(\ell_n\) can be constructed recursively.
\end{proof}

\subsection{Step 3. Construct the External Conformal Map}

\begin{theorem}[The External Conformal Map and Green Function]
\label{thm:external-conformal-map}
\(\M\) is nontrivial and connected, and its complement is connected.
In the category \(\mathbf{Riem}\) of Riemann surfaces and
biholomorphisms, there are mutually inverse isomorphisms
\[
 \Phi:\C\setminus\M\longrightarrow\{w\in\C:|w|>1\},
 \qquad \Psi:=\Phi^{-1},\qquad
 \frac{\Phi(c)}{c}\longrightarrow1\quad(|c|\to\infty).
\]
For \(k\in\N\), set
\[
 f_c^{\circ0}:=\id_{\C},\qquad
 f_c^{\circ(k+1)}:=f_c\circ f_c^{\circ k},\qquad
 g_c(z):=\lim_{k\to\infty}2^{-k}
             \log\max\{1,|f_c^{\circ k}(z)|\}.
\]
Define
\[
 G(c):=
 \begin{cases}
  \log|\Phi(c)|,&c\notin\M,\\
  0,&c\in\M.
 \end{cases}
\]
The functions \(g_c\) and \(G\) satisfy
 \begin{gather}
  G\in C(\C,[0,\infty)),\qquad G^{-1}(0)=\M,\qquad
  G(c)-\log|c|\longrightarrow0\quad(|c|\to\infty),\label{eq:green}\\
  g_c\circ f_c=2g_c,\qquad G(c)=g_c(c)=2g_c(0).
 \label{eq:green-normalization}
 \end{gather}
\end{theorem}
\begin{proof}
Connectedness of \(\M\) and the existence of \(\Phi\) follow from
\href{https://pi.math.cornell.edu/~hubbard/OrsayEnglish.pdf#page=64}
{Theorem~8.1, pp.~64--65} and
\href{https://pi.math.cornell.edu/~hubbard/OrsayEnglish.pdf#page=65}
{Corollary~8.3, p.~65}.\par
The extended dynamical Green function is continuous by
\href{https://pi.math.cornell.edu/~hubbard/OrsayEnglish.pdf#page=62}
{Proposition~8.1(b)} on page~62.\par
The normalization of the parameter Green function is given in
\href{https://pi.math.cornell.edu/~hubbard/OrsayEnglish.pdf#page=65}
{Corollary~8.4(a), p.~65}.
The functional equation follows by shifting the index in the defining
limit.
\end{proof}

\subsection{Step 4. Construct Retractions onto Green Sublevel Sets}

\begin{definition}[Green Level Curves and Retractions]\label{def:green-retractions}
For \(s>0\), set
\begin{gather}
 E_s:=\{c\in\C:G(c)\le s\},\qquad
 \kappa(s):=\dH(E_s,\M),\qquad
 \Gamma_s(\theta):=\Psi(e^{s+i\theta}),\label{eq:green-levels}\\
 P_s(c):=
 \begin{cases}
  c,&G(c)\le s,\\
  \displaystyle\Psi\left(e^s\frac{\Phi(c)}{|\Phi(c)|}\right),
       &G(c)>s,
 \end{cases}
 \qquad Q_s:=\Pi\circ P_s|_{\D}.
 \label{eq:green-projection}
\end{gather}
Let \(e_s:E_s\hookrightarrow\C\) be the inclusion. The map \(P_s\)
factors uniquely through a continuous map \(\bar P_s:\C\to E_s\) such
that \(e_s\circ\bar P_s=P_s\).
\end{definition}

\begin{theorem}[Retractions onto Green Sublevel Sets]
\label{thm:green-projection}
For \(0<s<t\), the following hold:
\begin{gather}
 E_s\subseteq\operatorname{int}E_t,\qquad
 \bigcap_{s>0}E_s=\M,\qquad
 \lim_{s\downarrow0}\kappa(s)=0,\label{eq:equipotential-distance}\\
 P_s(\C)=E_s,\qquad P_s|_{E_s}=\id_{E_s},\qquad
 \bar P_s\circ e_s=\id_{E_s},\qquad
 P_s\circ P_t=P_{\min\{s,t\}},\label{eq:retraction-composition}\\
 Q_s\in C(\D,\D),\qquad Q_s|_\M=\id_\M,\qquad
 \dH(Q_s(\D),\M)\le\kappa(s).\label{eq:Q-hausdorff-estimate}
\end{gather}
The boundary of \(E_s\) is the regular real-analytic Jordan curve
\(\Gamma_s\). The map \((s,c)\mapsto P_s(c)\) is continuous on
\((0,\infty)\times\C\). For every \(s>0\),
\[
 Q_s(\D)\setminus\M\ne\varnothing.
\]
For \(0<u<s<t\), the transition morphisms
\[
 r_{t,s}:=\bar P_s\circ e_t:E_t\to E_s
\]
obey \(r_{t,s}\circ e_{s,t}=\id_{E_s}\) and
\(r_{s,u}\circ r_{t,s}=r_{t,u}\), where \(e_{s,t}:E_s\hookrightarrow E_t\)
is the inclusion. Thus the maps \(r_{t,s}\) form a compatible system
of retractions.
\end{theorem}
\begin{proof}
Properness of \(G\) makes each \(E_s\) compact. The nesting follows
from the definition, and \(G^{-1}(0)=\M\) gives
\(\bigcap_{s>0}E_s=\M\). If \(\limsup_{s\downarrow0}\kappa(s)>0\),
choose \(s_j\downarrow0\) and \(c_j\in E_{s_j}\) with
\(\dist(c_j,\M)\ge\varepsilon\) for some \(\varepsilon>0\).
Compactness of \(E_{s_1}\) gives a convergent subsequence with limit
\(c_\infty\in G^{-1}(0)=\M\). Continuity of distance gives
\(\dist(c_\infty,\M)=0\), contradicting the lower bound
\(\dist(c_\infty,\M)\ge\varepsilon\).

The map \(\Gamma_s\) is injective on the circle, and
\[
 \Gamma_s'(\theta)=ie^{s+i\theta}\Psi'(e^{s+i\theta})\ne0.
\]
The unbounded component of the complement of this curve is
\(\Psi(\{|w|>e^s\})=\{G>s\}\). By the Jordan curve theorem, the filled
interior is \(E_s\).

For \(c\notin\M\), the map \(P_s\) can be written as
\[
 P_s(c)=
 \Psi\left(e^{\min\{s,G(c)\}}\frac{\Phi(c)}{|\Phi(c)|}\right).
\]
This proves continuity outside \(\M\), including across \(G=s\). At a
point \((s_0,c_0)\) with \(c_0\in\M\), continuity of \(G\) gives a
neighborhood on which \(G(c)<s_0/2\). For \(s\) near \(s_0\), \(P_s\)
is the identity on this neighborhood. The semigroup identity follows
by replacing the potential \(G(c)\) with \(\min\{s,G(c)\}\) while
leaving the external angle unchanged.

The map \(\Pi\) is a nonexpansive metric projection onto \(K\) and fixes
\(\M\subseteq K\). Hence
\[
 \dist(\Pi(z),\M)\le\dist(z,\M).
\]
Since \(P_s(\D)\subseteq E_s\) and \(Q_s\) fixes \(\M\), this proves
\eqref{eq:Q-hausdorff-estimate}.

Finally, \(c_0=\frac14\) is a boundary point of \(\M\), and there are
real parameters \(c_j>\frac14\) outside \(\M\) with \(c_j\to c_0\)
(as established in the boundary example of
Chapter~\ref{sec:remaining-factorization}). Since \(G(c_0)=0\) and
\(G\) is continuous, for every \(s>0\) some such \(c_j\) lies in
\(K\cap\{G<s\}\setminus\M\). For this point \(Q_s(c_j)=c_j\), so
\(c_j\in Q_s(\D)\setminus\M\).
\end{proof}

\subsection{Step 5. Embed \(O_N\) in Green Sublevel Sets}

\begin{lemma}[Embedding \(O_N\) in a Green Sublevel Set]
\label{lem:green-depth}
For every \(N\in\N\),
\begin{equation}\label{eq:green-depth-bound}
 \sup_{c\in O_N}G(c)\le2^{1-N}\log3.
\end{equation}
In particular,
\begin{equation}\label{eq:green-depth-condition}
 \begin{gathered}
 2^{1-N}\log3\le s\to
  \exists\eta_{N,s}:O_N\to E_s\text{ such that}\\
  e_s\circ\eta_{N,s}=j_N^\C,\qquad
  \bar P_s\circ j_N^\C=\eta_{N,s},\\
  Q_s|_{O_N}=\id_{O_N},\qquad h_N\le\kappa(s).
 \end{gathered}
\end{equation}
\noindent
Here \(j_N^\C:O_N\hookrightarrow\C\) denotes the inclusion into the
ambient space.
\end{lemma}
\begin{proof}
For \(|c|\le2\) and \(|z|\le2\), induction gives
\[
 |f_c^{\circ k}(z)|\le3^{2^k}-1.
\]
The base case is \(k=0\). If \(A=3^{2^k}\), the induction step follows
from \((A-1)^2+2\le A^2-1\). Thus \(g_c(z)\le\log3\) for these
\(c,z\). If \(c\in O_N\), then \(c,p_N(c)\in K\). By
\eqref{eq:green-normalization},
\[
 G(c)=2g_c(0)=2^{1-N}g_c(p_N(c))
                         \le2^{1-N}\log3.
\]
This gives the inclusion \(\eta_{N,s}:O_N\to E_s\).
The map \(\bar P_s\) satisfies
\(\bar P_s\circ j_N^\C=\eta_{N,s}\). Since \(O_N\subseteq K\),
\(\Pi\) fixes \(O_N\), so \(Q_s|_{O_N}=\id_{O_N}\). The inclusion
\(O_N\hookrightarrow E_s\) gives \(h_N\le\kappa(s)\).
\end{proof}

\subsection{Step 6. Construct a Sequence of Approximate Endomorphisms}

\begin{theorem}[Explicit External Morphisms]
\label{thm:constructed-external-maps}
Set
\begin{gather}
 a:=1+\max_{z\in\D}G(z),\qquad j_1:=0,\qquad s_1:=a,\notag\\
 j_n:=\min\{j\in\N:j>j_{n-1},\
                       \kappa(a2^{-j})\le2^{-n}\},\qquad
 s_n:=a2^{-j_n}\quad(n\ge2),\label{eq:height-choice}\\
 \lambda_0:=0,\qquad\lambda_1:=1,\notag\\
 \lambda_n:=\min\{N\in\N:N>\lambda_{n-1},\
                N\ge n,\ 2^{1-N}\log3\le s_n\}\quad(n\ge2),
 \label{eq:external-depths}\\
 T_0:=\id_{\D},\qquad T_1:=\Pi|_{\D}=Q_a,\qquad
 T_n:=Q_{s_n}\quad(n\ge2).\label{eq:external-maps}
\end{gather}
Then \(\lambda\in\Lambda\), each \(T_n\) is an endomorphism of
\(\D\) in \(\mathbf{CompHaus}\), and all assertions of the initial
theorem hold.
\end{theorem}
\begin{proof}
The minima defining \(j_n\) exist because \(\kappa(s)\to0\); hence
\(s_n\downarrow0\). The minima defining \(\lambda_n\) exist because
\(2^{1-N}\log3\to0\). Lemma~\ref{lem:green-depth} gives
\begin{equation}
 h_{\lambda_n}\le\kappa(s_n)\le2^{-n}\qquad(n\ge2).
 \label{eq:external-mesh-bound}
\end{equation}
The recursive definition and \eqref{eq:lower-depths} imply
\(\lambda_n\ge\ell_n\), so \(\lambda\in\Lambda\) and
\(\lambda_n\to\infty\). The sequence \((O_{\lambda_n})\) consists of
finite-orbit approximations.

Since \(G<a\) on \(\D\), we have \(P_a|_{\D}=\id_{\D}\) and
\(Q_a=\Pi|_{\D}\). At levels zero and one, \(O_0=O_1=K\), which are
fixed by the identity map and by \(\Pi\). For \(n\ge2\), the
factorization morphism \(\eta_n:O_{\lambda_n}\to E_{s_n}\) from
Lemma~\ref{lem:green-depth} satisfies
\(\bar P_n\circ\iota_{\lambda_n}=\eta_n\). Postcomposition with
\(\Pi\circ e_{s_n}\) preserves the identity
\(T_n\circ\iota_{\lambda_n}=\iota_{\lambda_n}\).
For \(z\in\D\)
\[
 \dist(T_n(z),O_{\lambda_n})
 \le\dist(T_n(z),\M)
 \le\kappa(s_n)\le2^{-n}.
\]
For \(n=0\), \(\dist(z,O_0)\le2\sqrt2-2<1\); for \(n=1\),
\(T_1(\D)=O_1=K\). Since every \(T_n\) fixes \(\M\), we have
\(\M\subseteq T_n(\D)\).
Finally, for \(n\ge2\),
\[
 \dH(T_n(\D),\M)\le\kappa(s_n)\le2^{-n}.
\]
\end{proof}

\begin{proposition}[Comparison of Finite Systems]
\label{prop:inverse-system-comparison}
For \(\lambda,s_n\) from Theorem~\ref{thm:constructed-external-maps},
set \(O_n^\lambda:=O_{\lambda_n}\) and \(E_n:=E_{s_n}\) for \(n\ge2\).
The inclusions define inverse functors
\[
 \mathcal O^\lambda,\mathcal E:
 \mathbf N_{\ge2}^{\mathrm{op}}\longrightarrow\mathbf{CompHaus}.
\]
The factorization morphisms \(\eta_n:O_n^\lambda\to E_n\) form a
natural transformation and induce the identity on the limits:
\[
 \iota_n^{\mathcal E}\circ\eta_{n+1}
       =\eta_n\circ\iota_n^{\mathcal O},\qquad
 \varprojlim\mathcal O^\lambda\cong\M
       \cong\varprojlim\mathcal E,\qquad
 \varprojlim\eta=\id_\M.
\]
If \(\bar P_n:=\bar P_{s_n}|_{\D}:\D\to E_n\), then
\[
 r_{n,m}:=\bar P_{s_m}|_{E_n}:E_n\to E_m\qquad(n\le m)
 \]
 define a covariant functor \(\mathcal S\), and \((\bar P_n)\) forms a
 natural transformation \(\underline{\D}\to\mathcal S\):
\[
 r_{n,n}=\id_{E_n},\qquad
 r_{m,k}\circ r_{n,m}=r_{n,k},\qquad
 r_{n,m}\circ\bar P_n=\bar P_m,\qquad
 r_{n,n+1}\circ\iota_n^{\mathcal E}=\id_{E_{n+1}}.
\]
For \(n\ge2\), the endomorphisms \(T_n\) from the approximation theorem
are given by
\(T_n=\Pi\circ e_{s_n}\circ\bar P_n\).
\end{proposition}
\begin{proof}
Lemma~\ref{lem:green-depth} gives inclusions
\(O_{\lambda_n}\hookrightarrow E_{s_n}\) on which \(\bar P_{s_n}\) is
the identity. Thus each \(\eta_n\) is an inclusion, and naturality
follows from commutativity of the inclusion diagrams. The stated limits
follow from \(\lambda_n\to\infty\) and \(s_n\downarrow0\). The induced
map \(\varprojlim\eta\) is the identity at every \(c\in\M\).

The remaining identities follow from
\eqref{eq:retraction-composition}: the transitions \(r_{n,m}\) are
compatible with composition and are left inverses to the inclusions
\(E_m\hookrightarrow E_n\); the identity
\(r_{n,m}\circ\bar P_n=\bar P_m\) expresses naturality of
\((\bar P_n)\). The final formula for \(T_n\) follows from
\eqref{eq:external-maps} and the definition of \(Q_{s_n}\).
\end{proof}

\subsection{Inner and Outer Cells}

\begin{equation}
 \label{eq:dyadic-grid-cells}
 \begin{gathered}
 n\ge2,\qquad J_n:=\{0,\ldots,2^n-1\},\qquad J_n^2:=J_n\times J_n,\\
 q_{n,a}:=
 \left[-2+\frac{4a}{2^n},-2+\frac{4(a+1)}{2^n}\right],\\
 Q_{n,(a,b)}:=
 \{x+iy:x\in q_{n,a},\ y\in q_{n,b}\}
 \qquad((a,b)\in J_n^2).
 \end{gathered}
\end{equation}
Each cell \(Q_{n,\alpha}\) is compact and satisfies
\begin{equation}
 \operatorname{diam}Q_{n,\alpha}\le8\cdot2^{-n}.
 \label{eq:dyadic-grid-diameter}
\end{equation}
Indeed, the differences of the real and imaginary parts of two points
in one cell are each at most \(4\cdot2^{-n}\), and
\(|x+iy|\le|x|+|y|\). Set
\[
 U_{n,\alpha}:=O_{\lambda_n}\cap E_{s_n}\cap Q_{n,\alpha},
 \qquad
 I_{n,\alpha}:=\M\cap Q_{n,\alpha}.
\]
Here \(U_{n,\alpha}\) is an outer cell in the common refinement of the
systems \((O_{\lambda_n})\) and \((E_{s_n})\), and \(I_{n,\alpha}\) is
its inner part. Define
\[
 \mathcal A_n^{\mathrm{out}}:=
   \{\alpha\in J_n^2:U_{n,\alpha}\ne\varnothing\},
 \qquad
 \mathcal A_n^{\mathrm{in}}:=
   \{\alpha\in J_n^2:I_{n,\alpha}\ne\varnothing\}.
\]
Each \(I_{n,\alpha}\) is compact, and
\(\mathcal A_n^{\mathrm{in}}\subseteq\mathcal A_n^{\mathrm{out}}\).
Let \(\mathbf C_n^{\mathrm{in}}\) and
\(\mathbf C_n^{\mathrm{out}}\) be the finite discrete categories with
object sets \(\mathcal A_n^{\mathrm{in}}\) and
\(\mathcal A_n^{\mathrm{out}}\), respectively. The parent map
\[
 p_n(a',b'):=\left(\left\lfloor\frac{a'}2\right\rfloor,
                       \left\lfloor\frac{b'}2\right\rfloor\right)
 \qquad((a',b')\in J_{n+1}^2)
\]
defines transition functors between levels: the inclusions
\(O_{\lambda_{n+1}}\subseteq O_{\lambda_n}\),
\(E_{s_{n+1}}\subseteq E_{s_n}\), and
\(Q_{n+1,\beta}\subseteq Q_{n,p_n(\beta)}\) induce inclusions
\(U_{n+1,\beta}\hookrightarrow U_{n,p_n(\beta)}\) and
\(I_{n+1,\beta}\hookrightarrow I_{n,p_n(\beta)}\).
All \(U_{n,\alpha}\) and \(I_{n,\alpha}\) are compact. The realization
functors \(\mathcal U_n:\mathbf C_n^{\mathrm{out}}\to\mathbf{CompHaus}\)
and \(\mathcal I_n:\mathbf C_n^{\mathrm{in}}\to\mathbf{CompHaus}\)
send \(\alpha\) to \(U_{n,\alpha}\) and \(I_{n,\alpha}\),
respectively. Inclusions between cells define natural transformations
\(\mathcal U_{n+1}\Rightarrow\mathcal U_n\circ p_n\) and
\(\mathcal I_{n+1}\Rightarrow\mathcal I_n\circ p_n\). The inclusions
\(I_{n,\alpha}\hookrightarrow U_{n,\alpha}\) define a natural
transformation \(\mathcal I_n\Rightarrow\mathcal U_n\circ j_n\), where
\(j_n:\mathbf C_n^{\mathrm{in}}\hookrightarrow\mathbf C_n^{\mathrm{out}}\);
these transformations are compatible with the transitions between
levels.

\begin{proposition}[Convergence of Inner and Outer Cells]
\label{prop:cellwise-inner-outer}
For every \(n\ge2\), the following hold:
\[
 \bigcup_{\alpha\in\mathcal A_n^{\mathrm{in}}}I_{n,\alpha}=\M,\qquad
 \bigcup_{\alpha\in\mathcal A_n^{\mathrm{out}}}U_{n,\alpha}=O_{\lambda_n}.
\]
For \(\alpha\in\mathcal A_n^{\mathrm{in}}\),
\[
 I_{n,\alpha}\subseteq U_{n,\alpha},
 \qquad
 \dH(I_{n,\alpha},U_{n,\alpha})\le8\cdot2^{-n}.
\]
For \(\alpha\in\mathcal A_n^{\mathrm{out}}\setminus
\mathcal A_n^{\mathrm{in}}\),
\[
 \sup_{z\in U_{n,\alpha}}\dist(z,\M)
 \le h_{\lambda_n}\le\kappa(s_n)\le2^{-n}.
\]
In particular,
\[
 \dH\left(
   \bigcup_{\alpha\in\mathcal A_n^{\mathrm{in}}}I_{n,\alpha},
   \bigcup_{\alpha\in\mathcal A_n^{\mathrm{out}}}U_{n,\alpha}
 \right)=h_{\lambda_n}\longrightarrow0.
\]
\end{proposition}
\begin{proof}
The covering properties follow from the grid covering \(\D\) and the
inclusions \(\M\subseteq O_{\lambda_n}\subseteq E_{s_n}\). If
\(\alpha\in\mathcal A_n^{\mathrm{in}}\), then
\(I_{n,\alpha}\subseteq U_{n,\alpha}\).
For \(z\in U_{n,\alpha}\), choose \(w\in I_{n,\alpha}\). Both points
lie in \(Q_{n,\alpha}\), so
\(|z-w|\le\operatorname{diam}Q_{n,\alpha}\le8\cdot2^{-n}\). This gives
the Hausdorff bound.

If \(\alpha\in\mathcal A_n^{\mathrm{out}}\setminus
\mathcal A_n^{\mathrm{in}}\), then
\(U_{n,\alpha}\subseteq O_{\lambda_n}\), hence
\(\sup_{z\in U_{n,\alpha}}\dist(z,\M)\le h_{\lambda_n}\).
The bounds \(h_{\lambda_n}\le\kappa(s_n)\le2^{-n}\) follow from
Lemma~\ref{lem:green-depth} and the choice of \(s_n,\lambda_n\). The
unions of cells are \(\M\) and \(O_{\lambda_n}\), respectively, and
\(\dH(O_{\lambda_n},\M)=h_{\lambda_n}\).
\end{proof}

\begin{theorem}[Inverse Limit of Cell-Address Sets]
\label{thm:cell-category-limit}
Let
\[
 \mathfrak A^{\mathrm{out}}:=
 \varprojlim(\mathcal A_n^{\mathrm{out}},p_n),\qquad
 \mathfrak A^{\mathrm{in}}:=
 \varprojlim(\mathcal A_n^{\mathrm{in}},p_n),
\]
where the finite index sets carry the discrete topology. Then
\(\mathfrak A^{\mathrm{in}}=\mathfrak A^{\mathrm{out}}\). For every
address \(a=(\alpha_n)\), the intersection
\[
 \bigcap_{n\ge2}U_{n,\alpha_n}
\]
consists of a single point \(\pi(a)\in\M\). The map
\(\pi:\mathfrak A^{\mathrm{out}}\to\M\) is continuous and surjective;
identifying addresses with the same realization point gives a
homeomorphism
\[
 \mathfrak A^{\mathrm{out}}/{\sim}\ \cong\ \M,
 \qquad a\sim b\leftrightarrow\pi(a)=\pi(b).
\]
\end{theorem}
\begin{proof}
For every address, the sets \(U_{n,\alpha_n}\) are nonempty, compact,
and nested. The finite intersection property for compact sets makes
their intersection nonempty; the diameter bounds \(8\cdot2^{-n}\)
show that it consists of a single point. This point belongs to every
\(O_{\lambda_n}\), and \(\lambda_n\to\infty\) gives
\(\bigcap_nO_{\lambda_n}=\M\). Thus every outer address has its index
at each level in \(\mathcal A_n^{\mathrm{in}}\), so
\(\mathfrak A^{\mathrm{out}}\subseteq\mathfrak A^{\mathrm{in}}\).
The reverse inclusion follows from
\(I_{n,\alpha}\subseteq U_{n,\alpha}\).

For \(c=x+iy\in\M\), set
\[
 a_n:=\min\left\{\left\lfloor\frac{2^n(x+2)}4\right\rfloor,2^n-1\right\},
 \qquad
 b_n:=\min\left\{\left\lfloor\frac{2^n(y+2)}4\right\rfloor,2^n-1\right\}.
\]
Then \((a_n,b_n)\in\mathcal A_n^{\mathrm{in}}\), and halving the
indices gives \(p_n(a_{n+1},b_{n+1})=(a_n,b_n)\). Thus these indices
define an address realizing \(c\), so \(\pi\) is surjective. If two
addresses agree at level \(n\), their realization points lie in the
same cell \(Q_{n,\alpha}\), whose diameter is at most
\(8\cdot2^{-n}\); this proves continuity of \(\pi\). The space
\(\mathfrak A^{\mathrm{out}}\) is compact as an inverse limit of finite
discrete spaces. The continuous surjection \(\pi\) induces a continuous
bijection \(\mathfrak A^{\mathrm{out}}/{\sim}\to\M\); compactness of the
domain and Hausdorffness of \(\M\) make this a homeomorphism.
\end{proof}

Proposition~\ref{prop:cellwise-inner-outer} gives an explicit
Hausdorff-distance bound for each inner cell and its corresponding
outer cell. For maps \(G_j\in\mathcal R_{m_j}\), condition
\(\mathsf{EA}\) requires a common equicontinuity modulus; by
Proposition~\ref{prop:equicontinuous-selection}, this holds exactly
when \(\LC(\M)\) holds. The Hausdorff bounds describe the sets
\(I_{n,\alpha}\) and \(U_{n,\alpha}\); to obtain the diagonal condition
\eqref{eq:diagonal-finite-scale}, one must additionally choose
\(\delta_i>0\) and prove that, for each \(k\), the first \(k\) estimates
hold simultaneously for some \(G\in\mathcal R_m\) at arbitrarily
large \(m\).

\subsection{Step 7. Estimate Distances between the Maps \(Q_s\)}

\begin{lemma}[Uniform Distance Bound between \(Q_s\) and \(Q_t\)]
\label{lem:band-motion}
For \(s>0\), define
\[
 V(s):=\max_{\theta\in\R/2\pi\mathbb Z}
                     e^s|\Psi'(e^{s+i\theta})|.
\]
Then \(V\) is finite and continuous on \((0,\infty)\). The curve
\(c_\theta(s):=\Psi(e^{s+i\theta})\) satisfies
\begin{equation}\label{eq:external-ray-ode}
 \frac{dc_\theta}{ds}
 =\frac{\Phi(c_\theta)}{\Phi'(c_\theta)}
 =\frac{\nabla G(c_\theta)}{|\nabla G(c_\theta)|^2},
 \qquad |\nabla G(c_\theta)|
       =\frac{|\Phi'(c_\theta)|}{|\Phi(c_\theta)|}>0.
\end{equation}
For \(0<s<t\), we have
\begin{equation}\label{eq:finite-band-motion}
 \|Q_t-Q_s\|_\infty
 \le\|\bar P_t|_{\D}-\bar P_s|_{\D}\|_{\infty,\C}
 \le\int_s^t V(u)\,du.
\end{equation}
\end{lemma}
\begin{proof}
The maximum in the definition of \(V(s)\) is taken over a compact
circle on which \(\Psi'\) is continuous and nonzero. Differentiating
\(\Phi(c_\theta(s))=e^{s+i\theta}\) gives the first equality in
\eqref{eq:external-ray-ode}. The gradient formula follows from
\(G=\log|\Phi|\) and the Cauchy--Riemann equations. If \(G(c)\le s\),
both maps fix \(c\). Otherwise their values lie on the same external
ray between levels \(s\) and \(\min\{t,G(c)\}\). Integrating the speed
along the ray gives the second estimate in
\eqref{eq:finite-band-motion}; nonexpansiveness of \(\Pi\) gives the
first.
\end{proof}

\begin{proposition}[Distances between Consecutive Maps \(Q_{s_n}\)]
\label{prop:consecutive-retraction-distances}
For the heights from Theorem~\ref{thm:constructed-external-maps},
with \(s_1=a\) and \(s_n\downarrow0\), define for \(n\ge2\)
\[
 d_n:=
 \|\bar P_n-\iota_n^{\mathcal E}\circ\bar P_{n+1}\|_{\infty,\C}
 =\sup_{z\in\D}|\bar P_n(z)-\bar P_{n+1}(z)|.
\]
Then for \(n\ge2\), \(N\ge2\), and \(N\ge1\), respectively,
\begin{align}
 d_n&\le\int_{s_{n+1}}^{s_n}V(u)\,du,\label{eq:retraction-distance-bound}\\
 \sum_{n=2}^{N}d_n
   &\le\int_{s_{N+1}}^{s_2}V(u)\,du,\label{eq:retraction-distance-sum}\\
 \sum_{n=1}^{N}\|Q_{s_{n+1}}-Q_{s_n}\|_\infty
   &\le\int_{s_{N+1}}^aV(u)\,du.
   \label{eq:retraction-motion-sum}
\end{align}
In particular, if \(\int_0^aV(u)\,du<\infty\), then both infinite sums
\(\sum_{n\ge2}d_n\) and
\(\sum_{n\ge1}\|Q_{s_{n+1}}-Q_{s_n}\|_\infty\) are finite.
\end{proposition}
\begin{proof}
For \(t=s_n\) and \(s=s_{n+1}\), estimate
\eqref{eq:finite-band-motion} gives the first bound. Summing over the
disjoint intervals \([s_{n+1},s_n]\) gives the second. Since \(s_1=a\),
the same summation for distances between \(Q_{s_n}\) gives the third.
The final assertion follows from convergence of the improper integral.
\end{proof}

\subsection{Step 8. Construct an Admissible Sequence for Each \(m\)}

\begin{definition}[Admissible Finite Sequences of Maps]
\label{def:finite-sequences}
Define
\[
 \Delta:=
 \{\delta\in(0,\infty)^\N:
   \forall k,\ \delta_k\le2^{-k},\
   \delta_{k+1}\le\delta_k/2\}.
\]
For \(\delta\in\Delta\), set
\[
 \mathcal K_\delta:=
 \left\{F\in C(\D,\D):
   \forall k\ \forall z,w\in\D,\
   |z-w|\le\delta_k\to|F(z)-F(w)|\le2^{-k}
 \right\}.
\]
Fix \(\lambda\in\Lambda\) and \(\delta\in\Delta\), and set
\[
 e_i(F):=\sup_{z\in\D}\dist(F(z),O_{\lambda_i}),\qquad
 \epsilon_i:=2^{-i},\qquad\beta_i:=8\,2^{-i}.
\]
For \(m\in\N\), define
\begin{align}
 \mathcal P_m(\lambda,\delta):=
 \left\{(F_0,\ldots,F_m)\in\mathcal K_\delta^{m+1}:
 \begin{gathered}
  F_0=\id_{\D},\quad
  (m\ge1\to F_1=\Pi|_{\D}),\\
  \forall i\in\{1,\ldots,m\},\quad F_i(\D)\subseteq K,\\
  \forall i\in\{0,\ldots,m\}\ \forall c\in O_{\lambda_i},\
                             F_i(c)=c,\\
  \forall i\in\{0,\ldots,m\},\quad e_i(F_i)\le\epsilon_i,\\
  \forall i<m,\quad\|F_{i+1}-F_i\|_\infty\le\beta_i
 \end{gathered}\right\},\label{eq:finite-sequences}\\
 \mathsf{UG}_m(\lambda,\delta):={}&
                  [\mathcal P_m(\lambda,\delta)\ne\varnothing].
 \label{eq:ug-m}
\end{align}
Compactness of \(\mathcal K_\delta\) and \(\mathcal P_m(\lambda,\delta)\)
follows from the next lemma.
\end{definition}

\begin{lemma}[Compactness of the Spaces of Admissible Finite Sequences]
\label{lem:compact-sequences}
\(\mathcal K_\delta\) is compact in the uniform topology, and each
\(\mathcal P_m(\lambda,\delta)\) is compact in the finite product
topology.
\end{lemma}
\begin{proof}
The conditions on \(\delta_k\) and \(F\) give a common modulus of
continuity for \(\mathcal K_\delta\); all values lie in the compact
space \(\D\). By the Arzelà--Ascoli theorem, the family is relatively
compact. Uniform limits preserve the defining conditions, so
\(\mathcal K_\delta\) is closed and compact.

For \(F,G\in C(\D,\D)\),
\[
 |e_i(F)-e_i(G)|\le\|F-G\|_\infty.
\]
The fixing, range, approximation-error, and successive-distance
conditions define closed subsets of the compact finite product
\(\mathcal K_\delta^{m+1}\). Hence \(\mathcal P_m(\lambda,\delta)\) is
compact.
\end{proof}

\begin{proposition}[An Admissible Sequence Exists for Every \(m\)]
\label{prop:finite-sequences-exist}
\[
 \forall m\in\N\ \exists\lambda^{[m]}\in\Lambda\
 \exists\delta^{[m]}\in\Delta,\qquad
 \mathsf{UG}_m(\lambda^{[m]},\delta^{[m]}).
\]
\end{proposition}
\begin{proof}
For \(m=0\), take \(\lambda^{[0]}=\lambda\) from
Theorem~\ref{thm:constructed-external-maps} and \(F_0=\id_{\D}\). For
\(m=1\), take \(\lambda^{[1]}=\lambda\) and
\((F_0,F_1)=(\id_{\D},\Pi|_{\D})\). In both cases set
\(\delta_k^{[m]}=2^{-k}\). Then \(\delta^{[m]}\in\Delta\), and all maps
in the sequence belong to \(\mathcal K_{\delta^{[m]}}\). We have
\(e_0(F_0)\le2\sqrt2-2<1\), and for \(m=1\),
\(e_1(F_1)=0\) and \(\|F_1-F_0\|_\infty<1\le\beta_0\). Thus the
conditions in Definition~\ref{def:finite-sequences} hold.

Now let \(m\ge2\), and take \(s_m,\lambda_m\) from Theorem
\ref{thm:constructed-external-maps}. Set \(F_0=\id_{\D}\),
\(F_1=\Pi|_{\D}\), and \(F_i=Q_{s_m}\) for \(2\le i\le m\). Define
\(\lambda^{[m]}_0=0,\lambda^{[m]}_1=1,
\lambda^{[m]}_2=\lambda_m\), and
\[
 \lambda^{[m]}_i:=
 \max\{\lambda^{[m]}_{i-1}+1,\ell_i\}\quad(i\ge3).
\]
Since \(\lambda\in\Lambda\) and \(\ell_m\ge\ell_2\), we have
\(\lambda_m\ge\ell_m\ge\ell_2\) and \(\lambda_m\ge m\ge2\). The
recursion gives \(\lambda^{[m]}_i\ge\ell_i\) for \(i\ge3\) and strict
increase, so \(\lambda^{[m]}\in\Lambda\). For \(2\le i\le m\),
\(\lambda^{[m]}_i\ge\lambda_m\), and therefore
\[
 O_{\lambda^{[m]}_i}\subseteq O_{\lambda_m}\subseteq E_{s_m}.
\]
Thus \(Q_{s_m}\) fixes the required outer sets, and the approximation
error at level \(i\) is at most
\(\kappa(s_m)\le2^{-m}\le2^{-i}\). At level zero,
\(e_0(F_0)\le2\sqrt2-2<1\), while at level one \(e_1(F_1)=0\). The
initial maps fix \(O_0=O_1=K\), and
\[
 \|F_1-F_0\|_\infty<1\le\beta_0,\qquad
 \|F_2-F_1\|_\infty\le\operatorname{diam}(K)=4=\beta_1.
\]
All successive distances after the second map are zero. The finite
family \(\{\id_{\D},\Pi|_{\D},Q_{s_m}\}\) is uniformly continuous.
Recursively choose sufficiently small \(\delta_k^{[m]}>0\), with
\(\delta_k^{[m]}\le2^{-k}\) and
\(\delta_{k+1}^{[m]}\le\delta_k^{[m]}/2\); then
\(\delta^{[m]}\in\Delta\) and every \(F_i\) belongs to
\(\mathcal K_{\delta^{[m]}}\).
\end{proof}

\begin{definition}[Terminal Maps of Admissible Sequences]
\label{def:last-maps}
For \(m\ge1\), set
\[
 \mathcal T_m:=
 \left\{G\in C(\D,K):
  \exists\lambda\in\Lambda\ \exists\delta\in\Delta\
  \exists(F_0,\ldots,F_m)\in\mathcal P_m(\lambda,\delta),\
  G=F_m
 \right\}.
\]
\end{definition}

\begin{lemma}[Terminal Maps Belong to \(\mathcal R_m\)]
\label{lem:last-maps-in-R}
\[
 \forall m\ge1,\qquad \mathcal T_m\subseteq\mathcal R_m.
\]
\end{lemma}
\begin{proof}
Let \(G\in\mathcal T_m\). By definition, there are
\(\lambda\in\Lambda\), \(\delta\in\Delta\), and
\((F_0,\ldots,F_m)\in\mathcal P_m(\lambda,\delta)\) such that
\(G=F_m\). Set \(N=\lambda_m\). Then \(N\ge m\),
\(G|_{O_N}=\id_{O_N}\), and
\[
 G(\D)\subseteq K,\qquad
 \sup_{z\in\D}\dist(G(z),O_N)=e_m(F_m)\le2^{-m}.
\]
Therefore \(G\in\mathcal R_m\).
\end{proof}

\begin{definition}[Extension Functional for the Next Index]
\label{def:extension-functional}
Fix \(\lambda\in\Lambda\), \(\delta\in\Delta\), \(m\ge1\), and
\(\sigma=(F_0,\ldots,F_m)\in\mathcal P_m(\lambda,\delta)\). Set
\[
 \mathcal H_{m+1}:=
 \{F\in\mathcal K_\delta:F(\D)\subseteq K,\
                    F|_{O_{\lambda_{m+1}}}=\id\},
\]
\[
 J_m(F;\sigma):=
 \max\left\{
 \frac{e_{m+1}(F)}{2^{-(m+1)}},
 \frac{\|F-F_m\|_\infty}{8\,2^{-m}}
 \right\},\qquad
 q_m(\sigma):=\inf_{F\in\mathcal H_{m+1}}J_m(F;\sigma).
\]
\[
 \{F\in\mathcal H_{m+1}:J_m(F;\sigma)\le1\}
 =
 \{F\in\mathcal K_\delta:
   (F_0,\ldots,F_m,F)\in\mathcal P_{m+1}(\lambda,\delta)\}.
\]
\end{definition}

\begin{proposition}[A Minimum for the Extension Functional]
\label{prop:extension-minimum}
For every \(\sigma\in\mathcal P_m(\lambda,\delta)\), the minimum
\(q_m(\sigma)\) is attained, and
\[
 \exists F\in\mathcal K_\delta:
 (F_0,\ldots,F_m,F)\in\mathcal P_{m+1}(\lambda,\delta)
 \quad\leftrightarrow\quad q_m(\sigma)\le1.
\]
For any \(\sigma=(F_0,\ldots,F_m)\) and
\(\tau=(G_0,\ldots,G_m)\) in \(\mathcal P_m(\lambda,\delta)\),
\[
 |q_m(\sigma)-q_m(\tau)|
 \le\frac{\|F_m-G_m\|_\infty}{8\,2^{-m}},
 \qquad q_m(\sigma)\le2^{m+3}.
\]
\end{proposition}
\begin{proof}
\(\mathcal H_{m+1}\) is closed in the compact space
\(\mathcal K_\delta\). It is nonempty: \(\Pi|_{\D}\in\mathcal K_\delta\),
maps \(\D\) into \(K\), and fixes \(O_{\lambda_{m+1}}\subseteq K\).
The functional \(J_m(\,\cdot\,;\sigma)\) is continuous, so its minimum
is attained. By the definition of \(J_m\), the condition
\(J_m(F;\sigma)\le1\) is equivalent to the two estimates
\[
 e_{m+1}(F)\le2^{-(m+1)},\qquad
 \|F-F_m\|_\infty\le8\,2^{-m}.
\]
Together with \(F\in\mathcal H_{m+1}\), these are equivalent to
\((F_0,\ldots,F_m,F)\in\mathcal P_{m+1}(\lambda,\delta)\).

For fixed \(F\in\mathcal H_{m+1}\), the triangle inequality gives
\[
 \left|\|F-F_m\|_\infty-\|F-G_m\|_\infty\right|
 \le\|F_m-G_m\|_\infty.
\]
Taking the maximum preserves this Lipschitz bound. Taking minima over
the common compact set \(\mathcal H_{m+1}\) gives the stated estimate
for \(q_m\).

Finally, \(\Pi(\D)=K\) and \(\operatorname{diam}(K)=4\), so
\(e_{m+1}(\Pi)\le4\) and \(\|\Pi-F_m\|_\infty\le4\), since
\(F_m(\D)\subseteq K\). Therefore
\[
 q_m(\sigma)\le J_m(\Pi;\sigma)
 \le\max\{4\cdot2^{m+1},\,4/(8\,2^{-m})\}=2^{m+3}.
\]
\end{proof}

\section{Equicontinuous Approximation and MLC}

\begin{proposition}[Local Connectivity Is Equivalent to \(\mathsf{EA}\)]
\label{prop:equicontinuous-selection}
\[
 \LC(\M)\quad\leftrightarrow\quad\mathsf{EA},
\]
\[
 \forall\rho,(m_j),(G_j),\qquad
 \mathsf{EAdata}\bigl(\rho,(m_j),(G_j)\bigr)
 \to
 \forall\varepsilon>0,\qquad
 \omega_\M\!\left(\frac{\rho(\varepsilon/3)}2\right)
 \le\frac{2\varepsilon}3.
\]
\end{proposition}
\begin{proof}
Suppose \(\mathsf{EAdata}(\rho,(m_j),(G_j))\) holds. For each \(j\),
choose \(N_j\ge m_j\) as in the definition of \(\mathcal R_{m_j}\).
Then
\[
 G_j(\D)\subseteq K,\qquad
 G_j|_{O_{N_j}}=\id_{O_{N_j}},\qquad
 \sup_{z\in\D}\dist(G_j(z),O_{N_j})\le2^{-m_j}.
\]
By the Arzelà--Ascoli theorem, a subsequence
\(G_{j_\nu}\to R\in C(\D,K)\) converges uniformly. For fixed \(k\),
we have \(N_{j_\nu}\ge k\) for all sufficiently large \(\nu\), hence
\(O_{N_{j_\nu}}\subseteq O_k\), and
\[
 \dist(R(z),O_k)
 \le\|R-G_{j_\nu}\|_\infty+2^{-m_{j_\nu}}
 \longrightarrow0.
\]
Therefore \(R(\D)\subseteq\bigcap_kO_k=\M\). Since
\(\M\subseteq O_{N_j}\) and \(G_j|_{O_{N_j}}=\id\), we have
\(R|_\M=\id_\M\).

The uniform limit preserves the modulus estimates:
\[
 |z-w|\le\rho(\eta)\to|R(z)-R(w)|\le\eta
 \qquad(\eta>0,\ z,w\in\D).
\]
Set \(r_\varepsilon:=\rho(\varepsilon/3)/2\). For
\(x,y\in\M\) with \(|x-y|<r_\varepsilon\), the set
\[
 C:=R\bigl(\D\cap\overline B(x,r_\varepsilon)\bigr)
\]
is compact and connected, contains \(x,y\), and lies in \(\M\). For
each \(z\in\D\cap\overline B(x,r_\varepsilon)\),
\[
 |R(z)-x|=|R(z)-R(x)|\le\varepsilon/3.
\]
Hence \(\operatorname{diam}C\le2\varepsilon/3\), and
\[
 \omega_\M(r_\varepsilon)\le2\varepsilon/3<\varepsilon.
\]
Theorem~\ref{thm:small-connected-compacts} gives \(\LC(\M)\).

Now suppose \(\LC(\M)\). By Carathéodory's theorem,
\(\Psi=\Phi^{-1}\) extends continuously to \(\{|w|\ge1\}\); denote
the extension by \(\widehat\Psi\)
\cite[Theorem~16.6]{Milnor}. Define
\[
 R(c):=
 \begin{cases}
  c,&c\in\M,\\
  \widehat\Psi\!\left(\Phi(c)/|\Phi(c)|\right),&c\notin\M.
 \end{cases}
\]
For \(c\notin\M\), set \(\xi=\Phi(c)/|\Phi(c)|\). As
\(c\to\partial\M\), we have \(G(c)\to0\), and thus
\[
 |R(c)-c|
 =|\widehat\Psi(\xi)-\widehat\Psi(e^{G(c)}\xi)|
 \longrightarrow0.
\]
Thus \(R:\C\to\M\) is continuous and \(R|_\M=\id_\M\). On the annulus
\(A_a=\{w:1\le|w|\le e^a\}\), define
\[
 \omega_a(r):=
 \sup_{\substack{w,v\in A_a\\|w-v|\le r}}
 |\widehat\Psi(w)-\widehat\Psi(v)|.
\]
Then \(\omega_a(r)\to0\) as \(r\downarrow0\). For \(0<s\le a\) and
\(c\notin\M\), set
\(\xi=\Phi(c)/|\Phi(c)|\) and \(r=\exp(\min\{s,G(c)\})\). Then
\(P_s(c)=\widehat\Psi(r\xi)\), \(R(c)=\widehat\Psi(\xi)\), and
\(|r\xi-\xi|\le e^s-1\). The maps \(P_s\) and \(R\) agree on \(\M\).
Since \(R(\D)\subseteq\M\subseteq K\) and \(\Pi\) is 1-Lipschitz, for
\(z\in\D\)
\[
 |Q_s(z)-R(z)|
 =|\Pi(P_s(z))-\Pi(R(z))|
 \le|P_s(z)-R(z)|
 \le\omega_a(e^s-1).
\]
Therefore
\[
 \|Q_s-R|_\D\|_\infty
 \le\omega_a(e^s-1)\longrightarrow0
 \qquad(s\downarrow0).
\]

Set \(\lambda_0=0,\lambda_1=1,s_1=a\). For each \(n\ge2\), choose
\(s_n\) and \(\lambda_n\) so that
\[
 0<s_n<\min\{s_{n-1}/2,2^{-n}\},\qquad
 \kappa(s_n)\le2^{-n},\qquad
 \|Q_{s_n}-R|_\D\|_\infty\le2^{-n},
\]
\[
 \lambda_n>\lambda_{n-1},\qquad
 \lambda_n\ge\ell_n,\qquad
 2^{1-\lambda_n}\log3\le s_n.
\]
Then \(\lambda\in\Lambda\) and \(\lambda_n\ge n\) for \(n\ge1\).
Set
\[
 F_1=\Pi|_\D,\qquad F_n=Q_{s_n}\quad(n\ge2).
\]
Then \(F_n\to R|_\D\) uniformly, so \(\{F_n:n\ge1\}\) is
equicontinuous. Define
\[
 \Omega(r):=
 \sup_{n\ge1}\ \sup_{\substack{z,w\in\D\\|z-w|\le r}}
 |F_n(z)-F_n(w)|,\qquad 0<r\le1,
\]
\[
 r_\eta:=\sup\{r\in(0,1]:\Omega(r)\le\eta\},\qquad
 \rho(\eta):=\min\{\eta,r_\eta/2\}\quad(\eta>0).
\]
Equicontinuity gives \(r_\eta>0\). Monotonicity of \(\Omega\) and the
definition of the supremum imply
\[
 0<\rho(\eta)\le\eta,\qquad \Omega(\rho(\eta))\le\eta.
\]
Lemma~\ref{lem:green-depth} gives \(Q_{s_n}|_{O_{\lambda_n}}=\id\).
For \(n\ge2\),
\[
 \sup_{z\in\D}\dist(F_n(z),O_{\lambda_n})
 \le\sup_{z\in\D}\dist(Q_{s_n}(z),\M)
 \le\kappa(s_n)\le2^{-n}.
\]
Moreover, \(F_n(\D)\subseteq K\). The map \(F_1=\Pi|_\D\) fixes
\(O_1=K\), maps \(\D\) into \(K\), and has zero distance from \(O_1\).
Thus \(F_n\in\mathcal R_n\) for every \(n\ge1\). Taking \(m_j=j\) and
\(G_j=F_j\) gives \(\mathsf{EAdata}(\rho,(m_j),(G_j))\).
\end{proof}

\section{Other Approaches Considered}
\label{sec:other-approaches}

\noindent
The approximation-theorem proof includes the functional \(q_m\),
compactness of the spaces of admissible sequences, and a criterion
using equicontinuity of a sequence of maps \(G_j\in\mathcal R_{m_j}\)
with \(m_j\to\infty\). The alternative approaches and estimates below
may help strengthen the proof.


\begin{description}
\item[Retractions \(Q_s\).] A bound on the distance between the maps
\(Q_s\) at different levels has been proved. Finiteness of
\(\int_0^aV(u)\,du\) would imply their uniform convergence as
\(s\downarrow0\). This route therefore reduces to estimating that
integral.

\item[Extending a sequence.] Local modifications of maps and convex
interpolation give error and uniform-distance bounds when an
admissible extension exists; the minimum \(q_m\) characterizes exactly
whether the next map can be added. A uniform bound is needed to extend
the sequence at every \(m\). The current estimate
\(q_m(\sigma)\le2^{m+3}\) grows with \(m\), leaving this quantitative
step open.

\item[Summable distances.] We considered the condition that the uniform
distances between \(G_j\) and \(G_{j+1}\) are summable as
\(m_j\to\infty\):
\[
 \sum_{j\ge1}\|G_{j+1}-G_j\|_\infty<\infty.
\]
Equicontinuity already gives a uniformly convergent subsequence by
Arzelà--Ascoli; summability of the distances provides an additional
quantitative route to uniform convergence.

\item[Homotopy invariants.] We considered the connected components and
higher homotopy groups of \(\mathcal P_m\) as \(m\to\infty\). If, for
fixed \((\lambda,\delta)\), all the sets
\(\mathcal P_m\) are nonempty, compactness of the corresponding
nested sets of sequences gives a compatible infinite sequence. If
the parameters depend on \(m\), an inverse system and quantitative
estimates would be needed.

\item[Boundary decompositions and gluing.]
We considered covers of \(\partial\M\) by compact sets with
one-point intersections, their images, and limits of the glued
complexes. Finite unions of such sets have topological dimension
\(\le1\), by \(\dim_{\mathrm{top}}(\partial\M)\le1\)
\cite{HurewiczWallman} and the theorem on the dimension of closed
unions; moreover, \(\M\) contains \(B(0,1/4)\). This route requires
explicit maps defining the identifications and estimates showing that
the diameters of the glued sets tend to zero.
\end{description}
